\documentclass[tikz,border=8pt]{standalone}
\usepackage{ctex}
\usepackage{amsmath}
\usetikzlibrary{calc, arrows.meta, 3d}

\begin{document}
\begin{tikzpicture}[
  scale=1.0, thick,
  every node/.style={font=\small},
  % 等轴测投影向量
  x={(1cm,0cm)}, y={(0cm,1cm)}, z={(0.35cm,0.25cm)}
]

%% ============================================================
%% 长方体三视图 + 立体图  布局：4 列
%% 尺寸：宽 a=2，深 b=1.2，高 c=1.5
%% ============================================================
\def\a{2}   % 宽
\def\b{1.2} % 深
\def\c{1.5} % 高

%% ---------- 列 0：立体图（等轴测） ----------
\begin{scope}[xshift=0cm]
  % 底面
  \draw[black] (0,0,0) -- (\a,0,0) -- (\a,0,\b) -- (0,0,\b) -- cycle;
  % 顶面
  \draw[black] (0,\c,0) -- (\a,\c,0) -- (\a,\c,\b) -- (0,\c,\b) -- cycle;
  % 竖棱（前四条）
  \draw[black] (0,0,0) -- (0,\c,0);
  \draw[black] (\a,0,0) -- (\a,\c,0);
  \draw[black] (\a,0,\b) -- (\a,\c,\b);
  \draw[black] (0,0,\b) -- (0,\c,\b);
  % 隐藏棱（灰色虚线）
  \draw[gray, dashed] (0,0,0) -- (0,0,\b);
  \draw[gray, dashed] (0,0,\b) -- (\a,0,\b);
  \draw[gray, dashed] (0,0,\b) -- (0,\c,\b);

  % 顶点标签
  \node[left, font=\footnotesize] at (0,\c,0) {$D_1$};
  \node[right, font=\footnotesize] at (\a,\c,0) {$C_1$};
  \node[right, font=\footnotesize] at (\a,\c,\b) {$B_1$};
  \node[left, font=\footnotesize] at (0,\c,\b) {$A_1$};
  \node[left, font=\footnotesize] at (0,0,0) {$D$};
  \node[right, font=\footnotesize] at (\a,0,0) {$C$};
  \node[right, font=\footnotesize] at (\a,0,\b) {$B$};
  %\node[left, font=\footnotesize] at (0,0,\b) {$A$};

  % 子图标题
  \node[font=\footnotesize, align=center] at (1.0, -0.5, 0) {立体图};
\end{scope}

%% ---------- 列 1：正视图（从正前方 z=大 看 xz 平面） ----------
\begin{scope}[xshift=4.0cm, yshift=0cm]
  % 正视图是宽 × 高 的矩形
  \draw[black] (0,0) rectangle (\a, \c);
  % 尺寸标注
  \draw[<->, gray, thin] (0, -0.35) -- (\a, -0.35) node[midway, below, font=\footnotesize] {宽 $a$};
  \draw[<->, gray, thin] (\a+0.3, 0) -- (\a+0.3, \c) node[midway, right, font=\footnotesize] {高 $c$};
  \node[font=\footnotesize, align=center] at (\a/2, -0.9) {正视图\\（从正面看）};
\end{scope}

%% ---------- 列 2：侧视图（从右侧看 yz 平面） ----------
\begin{scope}[xshift=7.5cm, yshift=0cm]
  % 侧视图是深 × 高 的矩形
  \draw[black] (0,0) rectangle (\b, \c);
  \draw[<->, gray, thin] (0, -0.35) -- (\b, -0.35) node[midway, below, font=\footnotesize] {深 $b$};
  \draw[<->, gray, thin] (\b+0.3, 0) -- (\b+0.3, \c) node[midway, right, font=\footnotesize] {高 $c$};
  \node[font=\footnotesize, align=center] at (\b/2, -0.9) {侧视图\\（从右侧看）};
\end{scope}

%% ---------- 列 3：俯视图（从上方看 xy 平面） ----------
\begin{scope}[xshift=10.2cm, yshift=0cm]
  % 俯视图是宽 × 深 的矩形
  \draw[black] (0,0) rectangle (\a, \b);
  \draw[<->, gray, thin] (0, -0.35) -- (\a, -0.35) node[midway, below, font=\footnotesize] {宽 $a$};
  \draw[<->, gray, thin] (\a+0.3, 0) -- (\a+0.3, \b) node[midway, right, font=\footnotesize] {深 $b$};
  \node[font=\footnotesize, align=center] at (\a/2, -0.9) {俯视图\\（从上方看）};
\end{scope}

%% ---------- 三投影规律提示 ----------
\begin{scope}[xshift=0cm, yshift=-2.4cm]
  \node[font=\footnotesize, align=left] at (6.5, 0)
    {三投影规律：正视图与侧视图\textbf{高平齐} \quad 正视图与俯视图\textbf{长对正} \quad 侧视图与俯视图\textbf{宽相等}};
\end{scope}

\end{tikzpicture}
\end{document}
